\chapter*{Introduction générale}
\addstarredchapter{Introduction générale}
\markboth{INTRODUCTION GÉNÉRALE}{INTRODUCTION GÉNÉRALE}

% TODO: situer ce rapport en 1 à 2 paragraphes -- contexte du stage/projet
% (formation universitaire, importance de cette étape).
[Présentez le contexte général : par exemple, la transition entre la formation universitaire et le monde professionnel,
la formation et l'établissement auxquels se rattache ce rapport, ainsi que le lieu du stage ou du projet.]

% TODO: décrire en 1 à 2 paragraphes la problématique traitée par l'entreprise
% d'accueil ou le projet, ainsi que son importance.
[Décrivez la problématique à laquelle répond le projet et les motivations associées.]

% TODO: présenter en 1 paragraphe la solution, le produit ou la plateforme,
% et préciser votre contribution ainsi que le périmètre de votre travail.
[Présentez le produit ou la plateforme, ainsi que votre rôle et le périmètre de votre travail.]

Ce rapport présente les principales étapes de mon travail et s'articule autour des chapitres suivants :

\begin{itemize}
    \item Le premier chapitre, intitulé \textbf{[Présentation du contexte]}, présente le contexte
    général : l'entreprise ou l'organisme d'accueil, le projet ou la plateforme,
    ainsi que le cadre et les objectifs du travail.

    \item Le deuxième chapitre, intitulé \textbf{[Environnement de développement et méthodologie]},
    décrit l'écosystème technique dans lequel le projet a été développé : le matériel,
    les outils collaboratifs et les technologies logicielles utilisées.

    \item Le troisième chapitre, intitulé \textbf{[Analyse des besoins et conception du système]},
    présente les besoins fonctionnels et non fonctionnels, le carnet de produit
    et l'architecture du système.

    \item Le quatrième et dernier chapitre, intitulé \textbf{[Réalisation et fonctionnalités]}, détaille les
    modules et les fonctionnalités réalisés, les choix de conception et les difficultés
    rencontrées, à l'aide de captures d'écran.

\end{itemize}
